\subsection{REDUCED ROOT SYSTEMS}

A root system is said to be reduced if every root of the system is indivisible (no. 3).

PROPOSITION 12. Assume that R is irreducible and reduced.

\begin{enumerate}
\def\labelenumi{(\roman{enumi})}
\item
  The ratio \(\frac{(\beta|\beta)}{(\alpha|\alpha)}\) for \(\alpha \in \mathbb{R}, \beta \in \mathbb{R}\) must take one of the values \(1, 2, \frac{1}{2}, 3, \frac{1}{3}\) .
\item
  The set of the \((\alpha |\alpha)\) for \(\alpha \in \mathrm{R}\) has at most two elements.
\end{enumerate}

Since R is irreducible, the transforms of a root by W(R) generate V (no. 2, Cor. of Prop. 5). Hence, for any roots \(\alpha, \beta\) , there exists a root \(\beta'\) such that \((\alpha|\beta') \neq 0\) and \((\beta'|\beta') = (\beta|\beta)\) . By formula (9) of no. 1 and the list of no. 3, \(\frac{(\beta'|\beta')}{(\alpha|\alpha)}\) takes one of the values \(1, 2, \frac{1}{2}, 3, \frac{1}{3}\) (recall that the system is assumed to be reduced). By multiplying \((x|y)\) by a suitable scalar, we can assume that \((\alpha|\alpha) = 1\) for certain roots and that the other possible values of \((\beta|\beta)\) for \(\beta \in \mathbb{R}\) are 2 and 3. The values 2 and 3 cannot both be attained, since in that case there would exist \(\beta \in \mathbb{R}, \gamma \in \mathbb{R}\) such that \(\frac{(\gamma|\gamma)}{(\beta|\beta)} = \frac{3}{2}\) , contrary to what we have seen above.

PROPOSITION 13. Assume that R is irreducible, non-reduced and of rank \(\geqslant 2\) .

\begin{enumerate}
\def\labelenumi{(\roman{enumi})}
\item
  The set \(\mathbf{R}_0\) of indivisible roots is a root system in V; this system is irreducible and reduced; and \(\mathrm{W}(\mathrm{R}_0) = \mathrm{W}(\mathrm{R})\) .
\item
  Let \(\mathbf{A}\) be the set of roots \(\alpha\) for which \((\alpha |\alpha)\) takes the smallest value \(\lambda\) . Then any two non-proportional elements of \(\mathbf{A}\) are orthogonal.
\item
  Let B be the set of \(\beta \in \mathbb{R}\) such that \((\beta |\beta) = 2\lambda\) . Then \(B \neq \emptyset\) , \(R_0 = A \cup B\) , \(R = A \cup B \cup 2A\) .
\end{enumerate}

If \(\alpha \in \mathbb{R} - \mathbb{R}_0\) , then \(\frac{1}{2}\alpha \in \mathbb{R}\) , but \(\frac{1}{2}\left(\frac{1}{2}\alpha\right) \notin \mathbb{R}\) (Prop. 8), so \(\frac{1}{2}\alpha \in \mathbb{R}_0\) . This proves that \(\mathbb{R}_0\) satisfies \((\mathrm{RS}_1)\) . It is clear that, for all \(\alpha \in \mathbb{R}\) , \(s_{\alpha, \alpha'}(\mathbb{R}_0) = \mathbb{R}_0\) , so \(\mathbb{R}_0\) satisfies \((\mathrm{RS}_{II})\) and \((\mathrm{RS}_{III})\) . Since \(\alpha \in \mathbb{R} - \mathbb{R}_0\) implies that \(\frac{1}{2}\alpha \in \mathbb{R}_0\) and since \(s_\alpha = s_{\alpha/2}\) , we have \(\mathrm{W}(\mathrm{R}) = \mathrm{W}(\mathrm{R}_0)\) . Thus \(\mathbb{R}_0\) is irreducible (Cor. of Prop. 5), and it is evidently reduced.

Since R is not reduced, there exists \(\alpha\in R_{0}\) such that \(2\alpha\in R\) . Since \(R_{0}\) is irreducible and \(\dim V\geqslant2\) , \(\alpha\) cannot be proportional or orthogonal to every root. Let \(\beta\in R_{0}\) be such that \(n(\beta,\alpha)\neq0\) and \(\beta\) is not proportional to \(\alpha\) . Changing \(\beta\) to \(-\beta\) if necessary, we can assume that \(n(\beta,\alpha)>0\) . Now \(\frac{1}{2}n(\beta,\alpha)=n(\beta,2\alpha)\in\mathbf{Z}\) , so \(n(\beta,\alpha)\in2\mathbf{Z}\) . From the list in no. 3, \(n(\beta,\alpha)=2\) , \((\beta|\beta)=2(\alpha|\alpha)\) . Since \(R_{0}\) is reduced, Prop. 12 shows that, for all \(\gamma\in R_{0}\) , either \((\gamma|\gamma)=(\alpha|\alpha)\) or \((\gamma|\gamma)=2(\alpha|\alpha)\) . Also, the above shows that, for all \(\gamma\in R-R_{0}\) , the vector \(\frac{1}{2}\gamma\) is an element of \(R_{0}\) such that \(\left(\frac{1}{2}\gamma|\frac{1}{2}\gamma\right)=(\alpha|\alpha)\) . Thus, \(\lambda=(\alpha|\alpha)\) , \(B\neq\varnothing\) , \(R_{0}=A\cup B\) , and \(R\subseteq A\cup B\cup2A\) ; on the other hand, if \(\gamma\in A\) , there exists \(g\in W(R)\) such that \(\gamma=g(\alpha)\) (Prop. 11), so \(2\gamma=g(2\alpha)\in R\) ; thus \(2A\subseteq R\) and \(R=A\cup B\cup2A\) . Finally, let \(\gamma,\gamma'\) be two non-proportional elements of A. Then

\[
n (2 \gamma , \gamma^ {\prime}) = 2 n (\gamma , \gamma^ {\prime}) = 4 n (\gamma , 2 \gamma^ {\prime}) \in 4 \mathbf {Z}, \text { and } | n (\gamma , \gamma^ {\prime}) | \leqslant 1
\]

since \(\gamma\) and \(\gamma'\) have the same length, so \(n(\gamma, \gamma') = 0\) and \((\gamma|\gamma') = 0\) .

PROPOSITION 14. Assume that R is irreducible and reduced, and that \((\alpha|\alpha)\) takes the values \(\lambda\) and \(2\lambda\) for \(\alpha\in R\) . Let A be the set of roots \(\alpha\) such that \((\alpha|\alpha)=\lambda\) . Assume that any two non-proportional elements of A are orthogonal. Then \(R_{1}=R\cup2A\) is an irreducible non-reduced root system and R is the set of indivisible roots of \(R_{1}\) .

It is clear that \(R_{1}\) satisfies \((RS_{I})\) and \((RS_{III})\) . We show that, if \(\alpha, \beta \in R_{1}\) , then \(\langle\alpha^{\sim}, \beta\rangle \in Z\) . This is clear if \(\alpha \in R\) . Since \((2\alpha)^{\sim} = \frac{1}{2}\alpha^{\sim}\) for \(\alpha \in A\) , it is also immediate if \(\alpha, \beta \in 2A\) . Finally, assume that \(\beta \in R\) and that \(\alpha = 2\gamma\) with \(\gamma \in A\) .

\begin{enumerate}
\def\labelenumi{\arabic{enumi})}
\item
  If \(\gamma = \pm \beta\) , then \(\langle \alpha^{\prime},\beta \rangle = \pm \frac{1}{2}\langle \gamma^{\prime},\gamma \rangle = \pm 1\) .
\item
  If \(\gamma\) is not proportional to \(\beta\) and if \(\beta \in A\) , the assumption on \(A\) implies that \(\langle \gamma^{\sim},\beta \rangle = 0\) , so \(\langle \alpha^{\sim},\beta \rangle = 0\) .
\item
  If \(\beta \in \mathbb{R} - \mathbb{A}\) , then \((\beta |\beta) = 2\lambda = 2(\gamma |\gamma)\) , so \(\langle \beta, \gamma^{\prime} \rangle\) is equal to 0, 2 or -2 by the list in no. 3. Thus \(\langle \beta, \alpha^{\prime} \rangle = \frac{1}{2}\langle \beta, \gamma^{\prime} \rangle \in \mathbf{Z}\) .
\end{enumerate}

Thus \(R_{1}\) is a root system in V, and the other assertions are clear.

